\chapter{Conclusion and Future Work}
\label{ch:conclusion}

\section{Summary}

This report presented \textsc{EdgeGuard}, a framework that combines signed Merkle-tree
attestation of quantized weights with an INT8 activation monitor and a rule-based threat
classifier, together with a testbed of attack simulators and language-model-shaped weight sets.
The evaluation, generated end to end from recorded runs, supports five conclusions.

\begin{enumerate}[leftmargin=*]
  \item \textbf{Cryptographic attestation does what it should, and localises exactly.} Every
        simulated weight modification was detected and attributed to exactly the tensors that
        changed, from single least-significant bits to perturbations spread over every tensor.
        This result does not depend on the simulation.
  \item \textbf{The monitor's false-positive rate is governed by the number of channels.} A
        per-channel $4\sigma$ rule over \MChannels{} channels flagged \EmpFPRZFour{} of clean
        inputs, matching the \PredFPRZFour{} predicted by \cref{eq:fwer}. A Bonferroni threshold
        reduced this to \EmpFPRBonf{} with no loss of detection in the simulated scenarios. For a
        real language model, with tens of thousands of channels, this correction is essential.
  \item \textbf{The two detectors are complementary.} Synergy inputs that kept activations near
        the clean distribution were flagged by the monitor at \MonSynergy{} and by the combined
        pipeline at \FlagSynergy{}.
  \item \textbf{The labels need better evidence, not better tuning.} Single-tensor supply-chain
        modifications cannot be distinguished from runtime flips without temporal evidence. The
        early/late rule labels pervasive input attacks as faults, and the entropy score cannot
        respond to shifted or narrowed activations.
  \item \textbf{Full attestation does not scale to per-inference use for LLMs.} At about
        \MerkThrMin{}--\MerkThrMax{}\,MiB/s, verifying 1.4\,GiB took 1.6\,s. The Merkle structure supports
        amortised verification with a deterministic bound of about 315 tokens at a 5\,ms-per-token
        budget, or probabilistic bounds under randomisation.
\end{enumerate}

Building the evaluation also uncovered three defects, two integer-width errors in the bit-flip
and cold-boot simulators and one algebraic error in the streaming covariance. Each had passed
tests that checked intended behaviour. All three were fixed and covered by regression tests
against independent references.

\section{Limitations}

The principal limitation is that activation-level results come from simulated activations whose
distributions were chosen together with the detector. They show that the pipeline functions as
designed, not that it detects real attacks. No real model was executed, no NPU was used, and no
adaptive attacker was evaluated. The memorisation score's premise, that extraction probes produce
distinctive late-layer activations, is a hypothesis that this project has not tested.

\section{Plan for the next phase}

\begin{figure}[t]
\centering
\begin{tikzpicture}[x=0.9cm, y=0.62cm]
\foreach \m/\lab in {0/Dec,1/Jan,2/Feb,3/Mar,4/Apr,5/May} {
  \node[font=\scriptsize, text=eggray] at (\m*2+1,0.7) {\lab};
  \draw[eggray!40] (\m*2,0.3) -- (\m*2,-8.4);
}
\draw[eggray!40] (12,0.3) -- (12,-8.4);
\foreach \y/\t/\s/\e/\c in {
  0/W1: label and explanation fixes/0/1/egaqua,
  1/W2: real-model activation capture (PyTorch)/0.5/3/egblue,
  2/W3: memorisation hypothesis test/2/4.5/egblue,
  3/W4: threshold calibration on real data/3/5/egblue,
  4/W5: amortised attestation + keyed MAC/3.5/6.5/egorange,
  5/W6: QNN adapter on Snapdragon device/5/9/egorange,
  6/W7: adaptive-attacker study/7/10/egblue,
  7/W8: thesis writing and submission/9.5/12/egaqua} {
  \node[font=\scriptsize, anchor=east] at (-0.1,-\y-0.6) {\t};
  \fill[\c, rounded corners=1pt] (\s,-\y-0.35) rectangle (\e,-\y-0.85);
}
\end{tikzpicture}
\caption{Indicative schedule for the next phase (December 2026 -- May 2027). Blue: evaluation on
real models; orange: systems work; aqua: fixes and writing.}
\label{fig:gantt}
\end{figure}

The next phase replaces simulation with measurement wherever possible. \Cref{fig:gantt} shows an
indicative schedule.

\begin{description}[leftmargin=*,style=nextline]
  \item[W1 — Fixes from this evaluation.] Replace the early-share rule with a comparison of
        anomaly rates; reword the signature line of explanations; bind dtype and leaf count into
        the hashes and add domain separation (\cref{rem:padding}); make Bonferroni the default
        threshold; emit structured evidence alongside labels.
  \item[W2 — Real activations.] Run a small open-weight model, quantized to INT8, under PyTorch
        with forward hooks. Capture activations from benign prompts for calibration and from
        published jailbreak and prompt-injection sets for evaluation. Measure the clean
        false-positive rate at the Bonferroni threshold on real, non-Gaussian activations.
  \item[W3 — Memorisation hypothesis.] Compare late-layer statistics during generation of known
        memorised sequences and matched controls (\cref{sec:disc:llm}). Keep, refit or drop the
        memorisation score according to the result.
  \item[W4 — Calibration on real data.] Fit thresholds with held-out empirical quantiles as well as
        Bonferroni, and compare diagonal, low-rank and full covariance models.
  \item[W5 — Amortised attestation.] Implement randomised, size-weighted partial verification with
        Merkle proofs and a keyed-MAC fast path. Measure detection latency against injected flips
        during generation.
  \item[W6 — Device integration.] Implement the QNN adapter against the Qualcomm AI Engine Direct
        SDK, expose selected intermediate tensors from the HTP backend, and measure overhead on a
        Snapdragon development device.
  \item[W7 — Adaptive attacker.] Optimise adversarial prompts against a model with the monitor's
        thresholds as a constraint, and report how much the monitor raises the attack cost.
  \item[W8 — Writing.] Consolidate the results into the final thesis and, if they warrant it, a
        workshop paper.
\end{description}

\section{Closing remark}

The most useful outcome of this phase is a clear map of where the project stands. The
cryptographic half of \textsc{EdgeGuard} is sound, and its main open problem is cost. The
statistical half is wired correctly and its false-positive behaviour is understood, but its
detection ability remains unproven until it meets real activations. The next phase is designed to
supply that evidence.
